% Preamble requirements (add to main.tex if not present):
%   \usepackage{tikz}
%   \usepackage{xcolor}
%   \usetikzlibrary{arrows.meta,shapes.rounded,positioning,fit}
%
\begin{figure*}[t]
\centering
\usetikzlibrary{arrows.meta,shapes.rounded,positioning,fit}
\begin{tikzpicture}[
  font=\small,
  >=Stealth,
  % Box styles
  inputbox/.style={
    draw, rounded corners=4pt,
    fill=gray!10, text width=3.6cm,
    align=left, inner sep=6pt,
    line width=0.7pt
  },
  pipelinebox/.style={
    draw, rounded corners=4pt,
    fill=blue!8, text width=2.4cm,
    align=center, inner sep=6pt,
    line width=0.7pt
  },
  huntbox/.style={
    draw, rounded corners=4pt,
    fill=orange!8, text width=3.4cm,
    align=left, inner sep=6pt,
    line width=0.7pt
  },
  % Level band styles — green→yellow→orange→red
  levelband/.style={
    draw, rounded corners=2pt,
    text width=12.8cm,
    align=left, inner sep=4pt 5pt,
    minimum height=0.62cm,
    line width=0.6pt
  },
  l1band/.style={levelband, fill=green!25,  draw=green!60!black},
  l2band/.style={levelband, fill=yellow!30, draw=yellow!70!black},
  l3band/.style={levelband, fill=orange!25, draw=orange!70!black},
  l4band/.style={levelband, fill=orange!50, draw=orange!80!black},
  l5band/.style={levelband, fill=red!20,   draw=red!60!black},
  % Column label style
  collabel/.style={font=\small\bfseries, align=center},
  % Annotations inside bands
  bandlabel/.style={font=\footnotesize\bfseries},
  bandmeta/.style={font=\footnotesize},
]

%% ─────────────────────────────────────────────
%% COLUMN LABELS (row 0)
%% ─────────────────────────────────────────────
\node[collabel] (lbl-in) at (0, 0)         {Shared Inputs};
\node[collabel] (lbl-pipe) at (5.35, 0)    {Independent Runs};
\node[collabel] (lbl-out) at (10.2, 0)     {Outputs};

%% ─────────────────────────────────────────────
%% LEFT COLUMN — Shared Inputs
%% ─────────────────────────────────────────────
\node[inputbox, below=0.35cm of lbl-in] (inputs) {
  \textbf{AoE2 World Data}\\[1pt]
  \hspace*{4pt}\textit{civs, units, techs,}\\
  \hspace*{4pt}\textit{maps, economy}\\[4pt]
  \textbf{SCOPE.md}\\[1pt]
  \hspace*{4pt}\textit{5+1 puzzles, casual}\\
  \hspace*{4pt}\textit{solo, The Monk}\\[4pt]
  \textbf{Toolkit Skills}\\[1pt]
  \hspace*{4pt}\textit{Stages 3–8}\\[4pt]
  \textbf{Profile Library}\\[1pt]
  \hspace*{4pt}\textit{29 v2 profiles}
};

%% ─────────────────────────────────────────────
%% MIDDLE COLUMN — Pipeline boxes
%% ─────────────────────────────────────────────
\node[pipelinebox] (pipe1) at (5.35, -1.05) {
  \textbf{Run 1}\\[3pt]
  Stage 3\\[1pt]
  \(\downarrow\)\\[1pt]
  Stage 4\\[1pt]
  \(\downarrow\)\\[1pt]
  Stage 8
};

\node[pipelinebox, below=0.5cm of pipe1] (pipe2) {
  \textbf{Run 2}\\[3pt]
  Stage 3\\[1pt]
  \(\downarrow\)\\[1pt]
  Stage 4\\[1pt]
  \(\downarrow\)\\[1pt]
  Stage 8
};

%% ─────────────────────────────────────────────
%% RIGHT COLUMN — Hunt outputs
%% ─────────────────────────────────────────────
\node[huntbox] (hunt1) at (10.2, -1.05) {
  \textbf{Hunt 1 ``Wololo''}\\[4pt]
  Answer: \textsc{wololo}\\[2pt]
  Meta: crossword\\
  \hspace*{10pt}grid extraction\\[2pt]
  Rounds: 1 (flat)\\[2pt]
  Feeders: \textsc{castle,}\\
  \hspace*{10pt}\textsc{onager, loom,}\\
  \hspace*{10pt}\textsc{tower, patrol}
};

\node[huntbox, below=0.5cm of hunt1] (hunt2) {
  \textbf{Hunt 2 ``Run 2''}\\[4pt]
  Answer: \textsc{grail}\\[2pt]
  Meta: index \(+\)\\
  \hspace*{10pt}anagram\\[2pt]
  Rounds: 3 (gated)\\[2pt]
  Feeders: \textsc{forge,}\\
  \hspace*{10pt}\textsc{march, siege,}\\
  \hspace*{10pt}\textsc{faith, relic}
};

%% ─────────────────────────────────────────────
%% ARROWS  inputs → pipelines
%% ─────────────────────────────────────────────
\draw[->, line width=0.8pt] (inputs.east) -- (pipe1.west);
\draw[->, line width=0.8pt] (inputs.east) -- (pipe2.west);

%% ─────────────────────────────────────────────
%% ARROWS  pipelines → hunts
%% ─────────────────────────────────────────────
\draw[->, line width=0.8pt] (pipe1.east) -- (hunt1.west);
\draw[->, line width=0.8pt] (pipe2.east) -- (hunt2.west);

%% ─────────────────────────────────────────────
%% FIVE-LEVEL COMPARISON HIERARCHY (below main diagram)
%% ─────────────────────────────────────────────

% Anchor: below the lower hunt box
\node[l1band, below=1.1cm of hunt2.south west, anchor=north west,
      xshift=-10.4cm] (l1) {%
  \bandlabel{Level~1 \;Structural}\quad
  \bandmeta{Hunt scale, narrator voice, format \textperiodcentered{} specification-driven \textperiodcentered{} stable across runs}%
};

\node[l2band, below=0.18cm of l1.south west, anchor=north west] (l2) {%
  \bandlabel{Level~2 \;Pool}\quad
  \bandmeta{Candidate pool composition, type distributions \textperiodcentered{} organizing principle variable (Jaccard\,=\,0.11)}%
};

\node[l3band, below=0.18cm of l2.south west, anchor=north west] (l3) {%
  \bandlabel{Level~3 \;Assignment}\quad
  \bandmeta{Answer word sets \textperiodcentered{} zero lexical overlap (Jaccard\,=\,0.00); partial domain adjacency (3/5 slots)}%
};

\node[l4band, below=0.18cm of l3.south west, anchor=north west] (l4) {%
  \bandlabel{Level~4 \;Mechanism}\quad
  \bandmeta{Extraction mechanism families \textperiodcentered{} 5 of 6 positions in distinct families; Jaccard\,=\,0.11}%
};

\node[l5band, below=0.18cm of l4.south west, anchor=north west] (l5) {%
  \bandlabel{Level~5 \;Micro-design}\quad
  \bandmeta{Clue formulations, visual grammar, flavor text \textperiodcentered{} systematically variable within every puzzle slot}%
};

%% Stability key (right of bands)
\node[font=\footnotesize, align=left, right=0.35cm of l1.east, anchor=west,
      yshift=0pt] (key1) {\textcolor{green!50!black}{\rule{8pt}{8pt}}\;Stable};
\node[font=\footnotesize, align=left, right=0.35cm of l3.east, anchor=west] (key3)
      {\textcolor{orange!70!black}{\rule{8pt}{8pt}}\;Variable};
\node[font=\footnotesize, align=left, right=0.35cm of l5.east, anchor=west] (key5)
      {\textcolor{red!60!black}{\rule{8pt}{8pt}}\;Most variable};

%% Downward arrow alongside bands (leftmost edge)
\draw[->, line width=1.0pt, gray!70]
  ([xshift=-0.45cm]l1.north west) -- ([xshift=-0.45cm]l5.south west)
  node[midway, left=2pt, font=\footnotesize\itshape, gray] {increasing\\[-2pt]variability};

\end{tikzpicture}

\caption{The same-input divergence paradigm. Shared inputs (left) feed two independent
pipeline runs (center), producing structurally distinct hunts (right). The five-level
comparison hierarchy (bottom) measures divergence from coarse architectural decisions
(Level~1) to individual clue formulations (Level~5). Color indicates observed stability:
green~= stable across runs, red~= variable.}
\label{fig:paradigm}
\end{figure*}
